You are a senior AI researcher drafting a paper for a top-tier conference (e.g., NeurIPS, ICML, CVPR, ICLR). 
Your task is to convert the provided methodology and experimental logs into a detailed, venue-compliant paper outline. You must output a single JSON object.

\vspace{1em}
\noindent\textbf{Your inputs are:}
\begin{enumerate}
    \item \texttt{idea.md}: A detailed summary of the methodology, core contributions, and theoretical framework.
    \item \texttt{experimental\_log.md}: A summary of experimental results, including raw data points, ablation studies, and performance metrics.
    \item \texttt{template.tex}: The target structure. You must use the section commands (e.g., \textbackslash{}section\{...\}) found here as your primary skeleton.
    \item \texttt{conference\_guidelines.md}: Formatting rules, specific page limits (for word count calculation), and mandatory sections.
\end{enumerate}

\vspace{0.5em}
\noindent\textbf{Processing Directives}
\vspace{0.5em}

\noindent\textbf{Global Instruction:} Do not analyze inputs in isolation. You must synthesize information across all provided documents for every step.

\vspace{1em}
\noindent\textbf{Directive 1: Plotting \& Visualization Plan}\\
\vspace{0.5em}
Synthesize \texttt{experimental\_log.md} and \texttt{idea.md} to identify the most compelling evidence.
\begin{itemize}
    \item Determine which figures are essential to visually prove the hypothesis (e.g., convergence rates, qualitative visual comparisons).
    \item The \texttt{plot\_type} MUST be exactly "plot" or "diagram". If it is a plot, specify the specific chart type (e.g., Radar Chart) inside the \texttt{objective}.
    \item The \texttt{data\_source} MUST be exactly "idea.md", "experimental\_log.md", or "both".
    \item Determine the ideal \texttt{aspect\_ratio} for each figure. The \texttt{aspect\_ratio} MUST be exactly one of: "1:1", "1:4", "2:3", "3:2", "3:4", "4:1", "4:3", "4:5", "5:4", "9:16", "16:9", "21:9".
    \item The \texttt{figure\_id} MUST be a semantically meaningful string identifier summarizing the plot contents, like "fig\_framework\_overview" or "fig\_ablation\_study\_parameter\_sensitivity". It MUST NOT contain the word "Figure".
    \item Output Focus: Create an array of objects for the \texttt{plotting\_plan} key.
\end{itemize}

\vspace{1em}
\noindent\textbf{Directive 2: Research Graph \& Investigation Strategy (Intro \& Related Work)}\\
\vspace{0.5em}
Provide search instructions for a downstream literature review agent to build a Research Graph. Do not write the actual paper content.

\vspace{0.5em}

\noindent Prevent Citation Overlap: Strictly separate the scope of the Introduction from Related Work to ensure the agent searches for different tiers of literature.
\begin{itemize}
    \item Introduction: Focuses on macro-level context (foundational papers, surveys).
    \item Related Work: Focuses on micro-level technical comparisons (recent SOTA baselines, benchmarks). 
\end{itemize}

\noindent Introduction Strategy (Macro-Level Context, 10-20 papers):
\begin{itemize}
    \item Hypotheses: Define the "Hook" (broad context) and "Problem Gap" to be verified. CRITICAL: Strictly scope the problem gap and claims to match the specific datasets and evaluations present in \texttt{experimental\_log.md}. Do not over-claim generalization.
    \item Search Directions: Provide 3-5 specific queries to find: 
    \begin{enumerate}
        \item Papers establishing the real-world impact or urgency of the problem gap.
        \item Good survey or review papers on the topic.
        \item 3-5 Foundational papers that established the sub-field.
    \end{enumerate}
\end{itemize}

\noindent Related Work Strategy (Micro-Level Technical Baselines, 30-50 papers):
\begin{itemize}
    \item Divide the field into 2-4 distinct methodology clusters that directly compete with or precede our approach.
    \item For each cluster, define:
    \begin{enumerate}
        \item Methodology Cluster Name: The technical category.
        \item SOTA Investigation: Instructions to find recent papers for conceptual context. CRITICAL TIMELINE RULE: Do not instruct searches for any papers published after \{cutoff\_date\}. Furthermore, do NOT instruct the search for new "competitors" to beat if they are not exclusively in \texttt{experimental\_log.md}.
        \item Limitation Hypothesis: The suspected failure point of these competing methods, based on \texttt{idea.md}.
        \item Limitation Search Queries: Highly specific, narrow queries to find papers documenting these exact limitations.
        \item The Bridge: How our proposed method resolves this specific limitation.
    \end{enumerate}
\end{itemize}
Output Focus: Populate the \texttt{intro\_related\_work\_plan} key.

\vspace{1em}
\noindent\textbf{Directive 3: Section Writing Plan \& Sizing Constraints}\\
\vspace{0.5em}
Outline the remaining sections (Abstract, Methodology, Experiments, Conclusion, Appendix) into a detailed structural plan.

\begin{itemize}
    \item Structural Hierarchy: If Subsection X.1 is created, X.2 is mandatory. Do not create orphaned subsections. Omit subsections entirely if a section does not require division.
    \item Content Specificity: Explicitly reference source materials. 
    \begin{itemize}
        \item \textit{Avoid:} "Describe the model."
        \item \textit{Require:} "Formalize the Temporal-Aware Attention mechanism using Eq. 3 from \texttt{idea.md}."
    \end{itemize}
    \item Mandatory Citations (\texttt{citation\_hints}): You must provide targeted citation hints for all external dependencies. Every hint must point to a single, unambiguous canonical paper. 
    \begin{itemize}
        \item Required Coverage (EXHAUSTIVE): You MUST explicitly create a targeted \texttt{citation\_hints} query for EVERY SINGLE dataset, optimizer, metric, and foundational architecture/model you mention, no matter how ubiquitous or obvious it seems (e.g., AdamW, ResNet, ImageNet, CLIP, Transformer, LLaMA, GPT, LLaVA). If it is in the \texttt{experimental\_log.md} or \texttt{idea.md}, it MUST have a citation hint.
        \begin{enumerate}
            \item All baseline methods compared against.
            \item All datasets evaluated on.
            \item All standard metrics utilized.
            \item All foundational algorithms, architectures (e.g., ResNet, Transformer), foundational models (e.g., LLMs, VLMs, Diffusion models), optimizers (e.g., AdamW), or frameworks built upon.
        \end{enumerate}
        \item Format Constraint \& Anti-Hallucination Rule: If you know the exact author and title, use "Author (Exact Paper Title)". DO NOT guess or hallucinate authors. If you do not know the exact author, use this format: "research paper or technical report introducing '[Exact Model/Dataset/Metric Name]'".
    \end{itemize}
    \item Output Focus: Populate the \texttt{section\_plan} key.
\end{itemize}

\noindent Guidelines on Scientific Depth \& Mathematical Rigor:
\begin{itemize}
    \item Grounded Formalization: Propose explicit subsections for rigorous mathematical formulations (e.g., loss functions, core algorithms, theoretical proofs). You must base these strictly on \texttt{idea.md} and \texttt{experimental\_log.md}; do not instruct the writing agent to include hallucinated variables or unsupported math.
\end{itemize}

\vspace{0.5em}
\noindent\textbf{Strict Output Format (JSON)}\\
You must output a single, valid JSON object with the following three top-level keys: "plotting\_plan", "intro\_related\_work\_plan", and "section\_plan".

\vspace{0.5em}
\noindent Example Output:

\begin{verbatim}
{
  "plotting_plan": [
    {
      "figure_id": "fig_teaser_fig_cross_modal_alignment_performance",
      "title": "Teaser: Cross-Modal Alignment Performance",
      "plot_type": "plot",
      "data_source": "experimental_log.md",
      "objective": "Visual summary (Radar Chart) demonstrating that our method 
                   achieves SOTA balance across 5 metrics.",
      "aspect_ratio": "16:9"
    }
  ],
  "intro_related_work_plan": {
    "introduction_strategy": {
      "hook_hypothesis": "Video-LLMs are currently the dominant paradigm for 
                         short clips.",
      "problem_gap_hypothesis": "Context window limits prevent scaling to >5s 
                                videos efficiently.",
      "search_directions": [
        "Find highly cited papers establishing the real-world impact of 
         context limits in video generation",
        "Search for published 'long-context video generation' surveys",
        "Identify foundational papers establishing causal video generation"
      ]
    },
    "related_work_strategy": {
      "overview": "Investigate three specific paradigms to build a graph 
                  proving the necessity of our Sliding-Window approach.",
      "subsections": [
        {
          "subsection_title": "2.1 Autoregressive Video Generation",
          "methodology_cluster": "Discrete Tokenization & Transformers",
          "sota_investigation_mission": "Identify the current SOTA 
                                        autoregressive models from 2024-2025. 
                                        Determine their maximum stable 
                                        generation length.",
          "limitation_hypothesis": "These models suffer from 'drift' or 'error 
                                   propagation' because they lack bidirectional 
                                   context.",
          "limitation_search_queries": [
            "Autoregressive video generation error propagation metrics",
            "Causal masking limitations in temporal video transformers"
          ],
          "bridge_to_our_method": "Our method introduces bidirectional blocks 
                                  to fix the hypothesized drift issue."
        },
        {
          "subsection_title": "2.2 Diffusion-Based Editing Frameworks",
          "methodology_cluster": "DDIM Inversion & Cross-Attention",
          "sota_investigation_mission": "Find recent papers using DDIM 
                                        inversion for editing. Identify the 
                                        standard benchmarks they use.",
          "limitation_hypothesis": "They fail at large structural changes 
                                   because cross-attention maps are too rigid.",
          "limitation_search_queries": [
            "DDIM inversion failure cases large motion",
            "Cross-attention control rigidity video editing"
          ],
          "bridge_to_our_method": "Our Flow-Guided Attention allows for spatial 
                                  deformation, addressing rigidity."
        }
      ]
    }
  },
  "section_plan": [
    {
      "section_title": "Abstract",
      "subsections": [
        {
          "subsection_title": "Abstract Content",
          "content_bullets": [
            "Briefly state the problem of temporal inconsistency.",
            "Introduce the proposed method.",
            "Highlight key results."
          ],
          "citation_hints": []
        }
      ]
    },
    {
      "section_title": "3. Methodology",
      "subsections": [
        {
          "subsection_title": "3.1 Temporal-Aware Attention Mechanism",
          "content_bullets": [
            "Define the query-key matching logic", 
            "Explain the masking strategy"
          ],
          "citation_hints": [
            "Vaswani et al. (Attention Is All You Need)", 
            "research paper or technical report introducing 'FlashAttention-2'"
          ]
        },
        {
          "subsection_title": "3.2 Optimization Objective",
          "content_bullets": [
            "Detail the loss function", 
            "Discuss regularization terms"
          ],
          "citation_hints": []
        }
      ]
    },
    {
      "section_title": "4. Experiments",
      "subsections": [
        {
           "subsection_title": "4.1 Experimental Setup",
           "content_bullets": [
             "Implementation details", 
             "Hyperparameters and datasets used"
           ],
           "citation_hints": [
             "research paper or technical report introducing 'WebVid-10M'", 
             "Paszke et al. (PyTorch: An Imperative Style, High-Performance 
              Deep Learning Library)",
             "research paper or technical report introducing 'AdamW optimizer'",
             "research paper or technical report introducing 'Jaccard Index'"
           ]
        },
        {
           "subsection_title": "4.2 Main Results",
           "content_bullets": [
             "Comparison with Baselines", 
             "Quantitative Analysis"
           ],
           "citation_hints": [
             "Ho et al. (Denoising Diffusion Probabilistic Models)",
             "research paper or technical report introducing 
             'AVSegFormer baseline'"
           ]
        }
      ]
    }
  ]
}
\end{verbatim}